\documentclass[11pt]{article}
\usepackage[margin=1in]{geometry}
\usepackage{amsmath,amssymb}
\usepackage{booktabs}
\usepackage{hyperref}

\title{A thalamic circuit represents dose-like responses induced by nicotine-related beliefs in human smokers}
\author{}
\date{}

\begin{document}
\maketitle

\begin{abstract}
Could non-pharmacological constructs such as beliefs impact brain activity in a dose-dependent manner like drugs do? Nicotine-addicted human smokers were instructed to believe an e-cigarette contained ``low'', ``medium'', or ``high'' nicotine while the actual nicotine content was kept constant. After vaping, participants performed a decision-making task during fMRI. Thalamic value-related activity increased parametrically with belief dose. Thalamus--ventromedial prefrontal cortex (vmPFC) functional coupling also scaled with belief dose. These findings demonstrate a dose-dependent relationship between a thalamic circuit and nicotine-related beliefs in humans, previously thought to apply only to pharmacological agents.
\end{abstract}

\section{Introduction}
Beliefs can alter behaviour and wellbeing and are represented by biological substrates in the brain \citep{r1,r2}. The placebo effect exemplifies how beliefs can improve outcomes without an active ingredient \citep{r5,r6,r7,r8,r9}, and beliefs have been shown to gate neural responses in human smokers \citep{r10,r11}. Pharmacological dose-response relationships are well established for drugs including nicotine \citep{r12}, alcohol \citep{r13}, and marijuana \citep{r14}. We tested whether nicotine-related beliefs can modulate brain activity in a dose-dependent manner. The thalamus contains one of the highest densities of nicotinic acetylcholine receptors (nAChRs) in the human brain \citep{r19,r20}, and PET imaging shows substantial nAChR occupancy even with denicotinized cigarettes \citep{r22,r23,r24}, suggesting cognitive constructs may drive addiction neurobiology.

\section{Methods}
\subsection{Design}
Twenty smokers (final $n=60$ scans across 20 smokers) were instructed before vaping that the e-cigarette contained ``low'', ``medium'', or ``high'' levels of nicotine, while actual nicotine content was fixed (1.2\%). After vaping for 20 minutes, participants performed a sequential investment fMRI task. A non-smoking healthy control (HC) group ($n=31$) performed the same task without vaping.

\subsection{Sanity checks}
Saliva samples for cotinine measurement (LC-MS/MS) and exhaled carbon monoxide (CO) were collected before and after vaping. Nicotine intake (mg) was computed as cartridge weight change times actual nicotine content (1.2\%). Subjective perceived nicotine strength was rated on a 10-point scale.

\subsection{fMRI task and analysis}
Sequential investment task with market return $r_t=(p_t - p_{t-1})/p_{t-1}$; participants made positive or negative bets. $|r_t|$ served as the parametric modulator of value. Whole-brain ANOVA with belief as main factor and $|r_t|$ as parametric modulator was thresholded at cluster-defining $p<0.005$ uncorrected, FWE cluster-corrected $p<0.05$. ROI analyses used independent anatomical masks for the thalamus \citep{r28}, nucleus accumbens, putamen, and caudate, plus a finer parcellation of the thalamus \citep{r29}. Permutation tests used 2,000 shuffled belief assignments; decoding used regularised linear discriminant analysis (rLDA) on thalamic multivoxel patterns, with 10,000 label permutations. Psychophysiological interaction (PPI) analysis with thalamus as seed investigated functional coupling. Control GLMs for visual and motor processing were constructed, and the analysis was repeated without spatial smoothing.

\section{Results}
\subsection{Instruction changed perceived nicotine strength without affecting intake or metabolism}
Perceived nicotine strength increased parametrically with instructed belief (low $=3.52 \pm 0.61$, medium $=4.52 \pm 0.41$, high $=5.82 \pm 0.47$; rmANOVA $F(2,38)=9.71$, $P=0.0004$, partial $\eta^2=0.34$, 90\% CI $0.12 \leq \eta \leq 0.48$). Nicotine intake: low $=0.928 \pm 0.56$\,mg, medium $=0.719 \pm 0.423$\,mg, high $=0.783 \pm 0.434$\,mg; rmANOVA $F(2,38)=1.806$, $P=0.178$. Cotinine change across conditions: rmANOVA $F(2,32)=0.959$, $P=0.393$. CO: rmANOVA $F(2,32)=0.364$, $P=0.698$.

\subsection{Dose-like response in the thalamus}
Whole-brain ANOVA found a parametric belief-dose effect on value-tracking activity in the thalamus (peak at MNI $x=-15$, $y=-19$, $z=-1$; $p<0.05$ FWE cluster-corrected at cluster-defining $p<0.005$, $k=50$, $P=0.006$). No other region showed the effect at the same threshold. ROI: low $=0.157 \pm 1.047$, medium $=0.601 \pm 0.714$, high $=2.914 \pm 0.865$ (rmANOVA $F(2,38)=3.62$, $P=0.036$, partial $\eta^2=0.16$, 90\% CI $0.0057 \leq \eta \leq 0.30$). Healthy controls (HC) $=2.318 \pm 4.258$; smokers-low vs.\ HC $t(49)=1.95$, $p=0.057$; smokers-medium vs.\ HC $t(49)=1.53$, $p=0.132$; smokers-high vs.\ HC $t(49)=-0.50$, $p=0.616$.

Permutation analysis of ROI beta estimates (2,000 permutations) showed true values ranked higher than the surrogate distribution ($P=0.002$). Thalamic sub-parcellation (FDR $q=0.05$): VPL, pulvinar, LGN, and CM all $P<0.05$.

rLDA decoded belief condition from thalamic multivoxel patterns at 49.3\% accuracy (chance 33.3\%, surrogate $33.1 \pm 6.3\%$, permutation $P=0.011$). Only VPL showed significant decoding after FDR (FDR $q=0.05$, $P=0.018$). Subjective perceived nicotine strength correlated with thalamic activity across all participants (Spearman $r=0.27$, $P=0.035$).

\begin{table}[h]
\centering
\caption{Perceived nicotine and thalamic activity by belief condition (mean $\pm$ SD).}
\begin{tabular}{lccc}
\toprule
Measure & Low & Medium & High \\
\midrule
Perceived strength (AU) & $3.52 \pm 0.61$ & $4.52 \pm 0.41$ & $5.82 \pm 0.47$ \\
Thalamus BOLD (AU) & $0.157 \pm 1.047$ & $0.601 \pm 0.714$ & $2.914 \pm 0.865$ \\
Nicotine intake (mg) & $0.928 \pm 0.56$ & $0.719 \pm 0.423$ & $0.783 \pm 0.434$ \\
\bottomrule
\end{tabular}
\end{table}

\subsection{Control analyses}
Button-press behaviour did not differ between belief conditions; thalamic motor and visual GLMs did not differ. Unsmoothed preprocessing preserved the belief effect.

\subsection{Striatum tracked value but not belief}
Ventral striatum tracked $|r_t|$ across all conditions ($P_{FDR} q<0.01$) but did not differentiate belief conditions (whole brain $P>0.05$). NAcc ROI: smokers-low $=1.228 \pm 3.329$, medium $=1.248 \pm 2.828$, high $=1.016 \pm 2.698$; HC $=1.781 \pm 2.138$; rmANOVA $F(2,38)=0.056$, $P=0.945$; permutation $P=0.94$. NAcc decoding accuracy 34.0\% (surrogate $32.1 \pm 6.3\%$, $P=0.372$). Putamen rmANOVA $F(2,38)=1.15$, $P=0.327$; caudate rmANOVA $F(2,38)=0.24$, $P=0.781$.

\subsection{Thalamus--vmPFC PPI scales with belief dose}
PPI with thalamus seed showed belief modulation of thalamus-vmPFC functional coupling at whole brain ($P_\text{SVC}<0.05$, FWE cluster-corrected at cluster-defining $p<0.005$, $P=0.041$) and ROI levels (peak at MNI $x=-11$, $y=50$, $z=-6$, $k=5$). A ventral striatum seed PPI did not yield significant coupling changes at the same threshold.

\section{Discussion}
Verbal instruction about nicotine strength produced a parametric dose-like modulation of thalamic activity and of thalamus--vmPFC functional coupling, resembling pharmacologically induced dose responses. The thalamus contains one of the highest densities of nAChRs \citep{r19,r20,r42,r43,r44,r45}, supporting its role as the primary substrate for belief-induced effects. Our findings extend prior observations that low or absent nicotine doses can still produce substantial nAChR occupancy \citep{r22,r23,r24}. The ventral striatum tracked reward but did not distinguish beliefs, likely because participants were not conditioned to e-cigarettes as in prior work with traditional cigarettes \citep{r10,r11}. The vmPFC--thalamus circuit encodes belief information in a parametric rather than binary manner, extending prior accounts of task-state representation in vmPFC \citep{r35,r36,r48} and its role in gating expectations \citep{r49}.

\bibliographystyle{plain}
\bibliography{references}

\end{document}
